\documentclass[UTF8]{ctexart}
\usepackage{xeCJK}
\usepackage{geometry}
\geometry{left=2cm,right=2cm,top=2.5cm,bottom=2.5cm}
\setlength{\headheight}{12.7pt}

\usepackage{titlesec}
\titleformat{\section}
  {\normalfont\large\bfseries}
  {\thesection}
  {1em}
  {}
\titlespacing*{\section}
  {0pt}
  {3.5ex plus 1ex minus .2ex}
  {2.3ex plus .2ex}

\usepackage{amsmath}
\usepackage{amssymb}
\usepackage{amsthm}
\usepackage{enumitem}
\setlist[enumerate]{itemsep=0pt,topsep=0pt,parsep=0pt,leftmargin=0pt,itemindent=2.5em}
\setlist[itemize]{itemsep=0pt,topsep=0pt,parsep=0pt,leftmargin=0pt,itemindent=2.5em}
\usepackage{fancyhdr}
\pagestyle{fancy}
\fancyhf{}
\usepackage{graphicx}
\usepackage{hyperref}
\usepackage{indentfirst}
\usepackage{lmodern}


\begin{document}
\setlength{\abovedisplayskip}{1pt}
\setlength{\belowdisplayskip}{1pt}

\section{可表函子}

\hypertarget{sec:定义1.1}{}
\noindent\textbf{定义1.1 (可表函子)}

给定范畴$\mathbf{C}$和其上的预层$F$. 如果存在$\mathbf{C}$中的对象$X$使得
\[
    h_X\cong F,
\]
就称$F$是\textbf{可表函子}, 称$X$是它的一个\textbf{表示对象}.

\vspace{10pt}

显然一个可表函子的不同表示对象在$\mathbf{C}$中同构.

\vspace{10pt}

\hypertarget{sec:定义1.2}{}
\noindent\textbf{定义1.2 (泛元素)}

给定范畴$\mathbf{C}$和其上的预层$F$.
对$\mathbf{C}$中任意的对象$X$和任意的$s\in F(X)$, 如果
\[
    H_{X,F}(s):h_X\cong F,
\]
就称$(X,s)$是$F$的\textbf{泛元素}.

\vspace{10pt}

\hypertarget{sec:定理1.1}{}
\noindent\textbf{定理1.1}

任给范畴$\mathbf{C}$和其上的预层$F$, $F$可表当且仅当它有泛元素.

\noindent{\kaishu 证明.}

\noindent{\kaishu 必要性.}

如果有$h_X\cong F$, 那么取一个自然同构$\theta:h_X\cong F$, 令$s=\theta_X(I_X)$.
则由\href{02 米田引理#nameddest=sec:定理2.1}{米田引理}得
\[
    H_{X,F}(s)=\theta:h_X\cong F,
\]
即$(X,s)$是$F$的泛元素.

\noindent{\kaishu 充分性.} 显然. \hfill $\qedsymbol$

\vspace{10pt}

\hypertarget{sec:定理1.2}{}
\noindent\textbf{定理1.2 (泛性质)}

给定范畴$\mathbf{C}$, 对象$X$, 预层$F$和$s\in F(X)$.
$(X,s)$是$F$的泛元素当且仅当它们具有\textbf{泛性质}:

对$\mathbf{C}$中任意的对象$A$和任意的$t\in F(A)$,
存在唯一的$\alpha\in\mathbf{C}(A,X)$使得$F(\alpha)(s)=t$.

\noindent{\kaishu 证明.}

\noindent{\kaishu 必要性.}

如果$(X,s)$是泛元素, 即$H_{X,F}(s):h_X\cong F$, 那么对$\mathbf{C}$中任意的对象$A$,
\[
    H_{X,F}(s)_A:\,\mathbf{C}(A,X)\to F(A),\quad\alpha\mapsto F(\alpha)(s)
\]
是双射. 从而对任意的$t\in F(A)$,
存在唯一的$\alpha\in\mathbf{C}(A,X)$使得$F(\alpha)(s)=t$.

\noindent{\kaishu 充分性.}

泛性质同样可以说明$H_{X,F}(s)_A$是双射, 即$H_{X,F}(s)$是自然同构,
$(X,s)$是泛元素. \hfill $\qedsymbol$

% 逻辑非常简洁, 但不知道这些是不是正确的思想.

\end{document}